\chapter{Toda brackets}
\label{ch:Toda}

The purpose of this chapter is to establish various Toda brackets
that are used elsewhere in this manuscript.
\rev{Tables \ref{tab:Toda} and \ref{tab:Toda-null} collect
all of this information in one place.}
Many Toda brackets can be easily computed from
the Moss Convergence Theorem \ref{thm:Moss}.
These are summarized in the tables without further
discussion.  
However, some brackets require more complicated arguments.
Those arguments are collected in this chapter.
\revv{
For easy reference, the lemmas in this chapter are labelled with degrees that match the degrees given in the tables.}

We will need the
following 
$\C$-motivic version of a 
theorem of Toda \cite{Toda62}*{Theorem 3.6} that applies to
symmetric Toda brackets.

\begin{thm}
\label{thm:Toda-symmetric}
Let $\alpha$ be an element of $\pi_{s,w}$, with $s$ even.
There exists an element $\alpha^*$ in $\pi_{2s+1,2w}$ such that
$\langle \alpha, \beta, \alpha \rangle$ contains
the product $\beta \alpha^*$ for all $\beta$ such that
$\alpha \beta$ \revv{equals zero}.
\end{thm}

\begin{cor}
\label{cor:2-symmetric}
If $2 \beta = 0$, then $\langle 2, \beta, 2 \rangle$
contains $\tau \eta \beta$.
\end{cor}

\begin{proof}
Apply Theorem \ref{thm:Toda-symmetric} to $\alpha = 2$.
We need to find the value of $\alpha^*$. 
Table \ref{tab:Massey} shows that
the Massey product $\langle h_0, h_1, h_0 \rangle$
equals $\tau h_1^2$.
The Moss Convergence Theorem \ref{thm:Moss} then shows that 
$\langle 2, \eta, 2 \rangle$ equals $\tau \eta^2$.
It follows that $\alpha^*$ equals $\tau \eta$.
\end{proof}


\begin{thm}
\label{thm:Toda}
Tables \ref{tab:Toda} \rev{and \ref{tab:Toda-null}} list 
some Toda brackets in the 
$\C$-motivic stable homotopy groups.
\end{thm}

\begin{proof}
The fourth column of the table gives information about the proof
of each Toda bracket.  

If the fourth column shows a Massey product, then the 
Toda bracket follows from the Moss Convergence Theorem \ref{thm:Moss}.
If the fourth column shows an Adams differential, then
the Toda bracket follows from the Moss Convergence Theorem \ref{thm:Moss},
using the mentioned differential.

A few Toda brackets are established elsewhere in the literature; 
specific citations are given in these cases.

Additional more difficult cases are established in the following lemmas.
\end{proof}

Tables \ref{tab:Toda} lists information about some Toda brackets
\rev{that do not contain zero,
while Table \ref{tab:Toda-null} lists information about some Toda brackets
that do contain zero.}
The third columns of the tables give elements of the
Adams $E_\infty$-page that detect elements of the Toda brackets.
The fourth columns of the tables give partial
information about indeterminacies, again by giving detecting elements
of the Adams $E_\infty$-page.  We have not completely analyzed the
indeterminacies of all brackets when the details are inconsequential
for our purposes.  The fifth columns indicate the proofs of 
the Toda brackets, and the sixth columns shows where each specific
Toda bracket is used in the manuscript.



\begin{lemma}
\label{lem:kappa,2,eta}
\revdeg{16, 9}
The Toda bracket $\langle \kappa, 2, \eta \rangle$ contains
zero, with indeterminacy generated by $\eta \rho_{15}$.
\end{lemma}

\begin{proof}
Using the Adams differential $d_3(h_0 h_4) = h_0 d_0$, 
the Moss Convergence Theorem \ref{thm:Moss} shows that the
Toda bracket is detected in filtration at least $3$.  The
only element in sufficiently high filtration is $P c_0$, 
which detects the product $\eta \rho_{15}$.  This product
lies in the indeterminacy, so the bracket must contain
zero.
\end{proof}

\begin{lemma}
\label{lem:kappa,2,eta,nu}
\revdeg{20, 11}
The Toda bracket $\langle \kappa, 2, \eta, \nu \rangle$
is detected by $\tau g$.
\end{lemma}

\begin{proof}
The subbracket $\langle 2, \eta, \nu \rangle$ is strictly zero,
since $\pi_{5,3}$ is zero.  
The subbracket 
$\langle \kappa, 2, \eta \rangle$ contains zero by Lemma 
\ref{lem:kappa,2,eta}.
Therefore, the fourfold bracket
$\langle \kappa, 2, \eta, \nu \rangle$
is well-defined.

Shuffle to obtain
\[
\langle \kappa, 2, \eta, \nu \rangle \eta^2 = 
\kappa \langle 2, \eta, \nu, \eta^2 \rangle.
\]
Table \ref{tab:Toda} shows that $\epsilon$ is contained
in the Toda bracket $\langle \eta^2, \nu, \eta, 2 \rangle$, so
the latter expression equals $\epsilon \kappa$,
which is detected by $c_0 d_0$.
It follows that 
$\langle \kappa, 2, \eta, \nu \rangle$
must be detected by $\tau g$.
\end{proof}

\begin{lemma}
\label{lem:nubar,sigma,2sigma}
\revdeg{23, 13}
The Toda bracket
$\langle \epsilon + \eta \sigma, \sigma, 2 \sigma \rangle$
contains zero, with indeterminacy generated
by $4 \nu \kappabar$ in $\{ P h_1 d_0 \}$.
\end{lemma}

\begin{proof}
Consider the shuffle
\[
\langle \epsilon + \eta \sigma, \sigma, 2 \sigma \rangle \eta =
(\epsilon + \eta \sigma) \langle \sigma, 2 \sigma, \eta \rangle.
\]
Table \ref{tab:Toda} shows that $h_1 h_4$ detects
$\langle \sigma, 2 \sigma, \eta \rangle$, so
$h_1 h_4 c_0$ detects the product
$\epsilon \langle \sigma, 2 \sigma, \eta \rangle$.
On the other hand, Table \ref{tab:misc-extn} shows that
\revv{there is a hidden $\sigma$ extension from $h_1 h_4$ to $h_4 c_0$.
Therefore,}
$h_1 h_4 c_0$ also detects 
$\eta \sigma \langle \sigma, 2 \sigma, \eta \rangle$.
It follows that
$(\epsilon + \eta \sigma) \langle \sigma, 2 \sigma, \eta \rangle$
is detected in filtration greater than $5$.

\revv{Consider the shuffle
\[
2 \langle \epsilon + \eta \sigma, \sigma, 2 \sigma \rangle =
\langle 2, \epsilon + \eta \sigma, \sigma \rangle 2 \sigma.
\]
The latter expression is zero since $2 \sigma$ annihilates
all elements of $\pi_{16,9}$.
}
This shows that no elements of the Toda bracket
can be detected by $\tau h_2 g$ or $\tau h_0 h_2 g$.

The element $4 \nu \kappabar$ generates the indeterminacy
because it equals $\tau \eta \kappa (\epsilon + \eta \sigma)$.
\end{proof}

\begin{lemma}
\label{lem:tkappa^2,eta,2}
\revdeg{30, 16}
The Toda bracket $\langle \tau \kappa^2, \eta, 2 \rangle$
equals zero, with no indeterminacy.
\end{lemma}

\begin{proof}
The Adams differential $d_3(\D h_2^2) = \tau h_1 d_0^2$
implies that the bracket is detected by $h_0 \cdot \D h_2^2$,
which equals zero \revv{in the $E_\infty$-page}.  
Therefore, the Toda bracket is detected 
in Adams filtration at least $7$, but there are no 
elements in the Adams $E_\infty$-page in sufficiently high
filtration.

The indeterminacy can be computed by inspection.
\end{proof}

\begin{lemma}
\label{lem:eta^2,theta4,eta^2}
\revdeg{35, 20}
The Toda bracket $\langle \eta^2, \theta_4, \eta^2 \rangle$
contains zero, with indeterminacy generated by
$\eta^3 \eta_5$.
\end{lemma}

\begin{proof}
If the bracket were detected by $h_2 d_1$, then
\[
\nu \langle \eta^2, \theta_4, \eta^2 \rangle =
\langle \nu, \eta^2, \theta_4 \rangle \eta^2
\]
would be detected by $h_2^2 d_1$.
However, $h_2^2 d_1$ does not detect a multiple of $\eta^2$.

The bracket cannot be detected by $\tau h_1 e_0^2$ by 
comparison to $\tmf$.

By inspection, the only remaining possibility is that
the bracket contains zero.  The indeterminacy can be computed
by inspection.
\end{proof}

\begin{lemma}
\label{lem:t,eta^2kappa1,eta}
\revdeg{36, 20}
The Toda bracket $\langle \tau, \eta^2 \kappa_1, \eta \rangle$
is detected by $t$, with indeterminacy generated by
$\eta^3 \mu_{33}$.
\end{lemma}

\begin{proof}
There is a relation
$h_1 \cdot \ol{h_1^2 d_1} = t$ in the homotopy of $C\tau$.
Using the connection between Toda brackets and cofibers
as described in \cite{Isaksen14c}*{Section 3.1.1},
this shows that $t$ detects the Toda bracket.

The indeterminacy is computed by inspection.
\end{proof}

\rev{
The third author presents the following Lemma~\ref{lem:theta4,2,sigma^2+kappa} as a correction to \cite{Xu16}*{Theorem 2.1}, where it states that the said Toda bracket contains 0.

\begin{lemma}
\label{lem:theta4,2,sigma^2+kappa}
\revdeg{45}
The classical Toda bracket $\langle \theta_4, 2, \sigma^2 + \kappa \rangle$
contains $0$ or $\eta \kappabar_2$.  Its indeterminacy
is generated by $\rho_{15} \theta_4$, which is detected by $h_0^2 h_5 d_0$.
\end{lemma}

\begin{proof}
The gap originated in \cite{Xu16}*{Remark~3.3},
where it was claimed that $\langle \theta_4, 2, \sigma^2 \rangle$ contains an order 2 element of the form $2 \alpha + \beta$, where $\alpha$ is detected by $h_4^3$ and $\beta$ is detected by $h_5 d_0$. In fact, since $h_5 d_0$ and $h_1 g_2$ are in the same filtration, we can only conclude that $\beta$ is detected by $h_5 d_0$ or $h_5 d_0 + h_1 g_2$, therefore the missed possibility in the statement of the lemma.
\end{proof}
}

\rev{
\begin{remark}
\label{rem:theta4,2,sigma^2+kappa}
In fact, we have evidence that this classical Toda bracket\\
$\langle \theta_4, 2, \sigma^2 + \kappa \rangle$ contains $\eta \kappabar_2$.  However, the argument depends on computations
as far as the 110-stem.
\end{remark}
}

\begin{lemma}
\label{lem:eta,2,4kappabar2}
\revdeg{46, 25}
The Toda bracket
$\langle \eta, 2, 4 \kappabar_2 \rangle$
contains zero.
\end{lemma}

\begin{proof}
The Massey product $M h_1 = \langle h_1, h_0, h_0^2 g_2 \rangle$
shows that $M h_1$ detects the Toda bracket.
Table \ref{tab:eta-extn} shows that
$M h_1$, $\D h_2 c_1$, and $\tau d_0 l + \D c_0 d_0$ are all
targets of hidden $\eta$ extensions.
(Beware that the
hidden $\eta$ extension from $h_3^2 h_5$ to $M h_1$ is a crossing
extension in the sense of Section \ref{sctn:assoc-graded},
but that does not matter.)
Therefore, $M h_1$ detects only multiples of $\eta$,
so the Toda bracket contains a multiple of $\eta$.
This implies that it contains zero, since multiples of $\eta$
belong to the indeterminacy.
\end{proof}

\begin{lemma}
\label{lem:tkappabar2,sigma^2,2}
\revdeg{59, 31}
The Toda bracket $\langle \tau \kappabar_2, \sigma^2, 2 \rangle$
equals zero.
\end{lemma}

\begin{proof}
No elements of the bracket can
be detected by $\tau^2 \D h_1 d_0 g$ by comparison to $\tmf$.

Consider the shuffle
\[
\langle \tau \kappabar_2, \sigma^2, 2 \rangle \kappa =
\tau \kappabar_2 \langle \sigma^2, 2, \kappa \rangle.
\]
The bracket $\langle \sigma^2, 2, \kappa \rangle$ is zero
because it is contained in $\pi_{29,16} = 0$.
On the other hand,
$\{\tau M d_0\} \kappa$ is non-zero and detected by
$\tau M d_0^2$.
Therefore, no elements of
$\langle \tau \kappabar_2, \sigma^2, 2 \rangle$
can be can be detected by $\tau M d_0$.
\end{proof}

\rev{
The third author presents the following Lemma~\ref{lem:etakappabar2,2sigma,sigma}, which is needed in the proof of Lemma~\ref{lem:theta4-h4^2}.

\begin{lemma}
\label{lem:etakappabar2,2sigma,sigma}
\revdeg{60}
The classical Toda bracket $\langle \eta \kappabar_2, 2 \sigma, \sigma \rangle$
equals zero.
\end{lemma}

\begin{proof}
Due to $d_3(e_1) = h_2^2 n$ and that $h_1g_2 = h_3e_1$, we have the following Massey product in the $E_4$-page
$$h_1g_2 = \langle h_2, h_2n, h_3 \rangle$$
with zero indeterminacy. Since there are no crossing differentials, we conclude that $h_1g_2$ detects a homotopy class in the Toda bracket $\langle \nu, \nu\{n\}, \sigma \rangle$. We claim that $\eta \kappabar_2$ is contained in this bracket. In fact, they might differ by classes detected in filtration 6 or higher: $h_0h_5d_0, \ h_0^2 h_5d_0, \ w$. The first two detect $\sigma$-multiples so they are in the indeterminacy. The homotopy class $\{w\}$ is detected by tmf, but $\kappabar_2$ and $\{n\}$ are not, so $w$ can be ruled out too. 

Therefore, we have 
$$
\langle \eta \kappabar_2, 2\sigma, \sigma \rangle \subseteq \langle \langle \nu, \nu\{n\}, \sigma \rangle, 2\sigma, \sigma \rangle
\supseteq \nu \langle \nu\{n\}, \sigma, 2\sigma, \sigma \rangle = 0.
$$
Here $\langle \sigma, 2\sigma, \sigma \rangle = 0$, and $ \langle \nu\{n\}, \sigma, 2\sigma \rangle$ contains 0 since coker J in $\pi_{49}$ is 0. So the 4-fold bracket $\langle \nu\{n\}, \sigma, 2\sigma, \sigma \rangle$ in $\pi_{57}$ is well-defined. By comparison with $\pi_{57}tmf$, we know it is 0.

We remain to show the indeterminacy of $\langle \langle \nu, \nu\{n\}, \sigma \rangle, 2\sigma, \sigma \rangle$ is 0. In fact,
\begin{itemize}
\item
$\pi_{15} \cdot \langle \nu, \nu \{n\}, \sigma \rangle = \langle \pi_{15},  \nu, \nu \{n\} \rangle \sigma  \subseteq \sigma \pi_{53} = 0$.  (Lemma~2.3 in \cite{Xu16}.)

\item $\langle \nu \cdot \pi_{42}, 2\sigma, \sigma \rangle = \langle 0, 2\sigma, \sigma \rangle + \langle \rho_{15}\theta_4, 2\sigma, \sigma \rangle = 0$.  (Lemmas~2.3 and 2.4 in \cite{Xu16}.)

\item $\langle \sigma \cdot \pi_{38}, 2\sigma, \sigma \rangle \supseteq \pi_{38} \cdot \langle \sigma, 2\sigma, \sigma \rangle = 0$.
\end{itemize}
This completes the proof.
\end{proof}
}

\begin{lemma}
\label{lem:2,sigma^2,theta4.5}
\revdeg{60, 32}
For every $\alpha$ that is detected by $h_3^2 h_5$,
the Toda bracket $\langle 2, \sigma^2, \alpha \rangle$
contains zero.  The indeterminacy is generated by 
$2 \tau \kappabar^3$, which is detected by $\tau^2 d_0^2 l$.
\end{lemma}

\begin{proof}
\revv{
Let $\alpha$ be detected by $h_3^2 h_5$.
For degree reasons,
the only elements that could detect $\sigma^2 \alpha$ either support
$\eta$ extensions or are detected by $\tmf$.
Therefore, $\sigma^2 \alpha$ is zero.
}
Hence the bracket is defined.

By comparison to $\tmf$, the bracket cannot be detected
by $\tau^4 g^3$. 
Table \ref{tab:2-extn} shows that
$\tau^2 d_0^2 l$ is the target of a hidden $2$ extension, so it
detects an element in the indeterminacy.
Since there are no other possibilities, the bracket must
contain zero.
\end{proof}

\begin{remark}
This result is consistent with Table 23 of \cite{Isaksen14c},
which claims that the bracket
$\langle 2, \sigma^2, \theta_{4.5} \rangle$
contains an element that is detected by $B_3$.
The element $B_3$ is now known to be zero in the Adams $E_\infty$-page,
so this just means that the bracket contains an element detected
in Adams filtration strictly greater than the filtration of $B_3$.
\end{remark}

\begin{lemma}
\label{lem:theta4,eta^2,theta4}
deg{63, 34}
The Toda bracket $\langle \theta_4, \eta^2, \theta_4 \rangle$
equals zero.
\end{lemma}

\begin{proof}
Theorem \ref{thm:Toda-symmetric} says that there exists an element
$\theta_4^*$ in $\pi_{61,32}$ such that
$\langle \theta_4, \eta^2, \theta_4 \rangle$ contains
$\eta^2 \theta_4^*$.
The group $\pi_{61,32}$ is zero, so $\theta_4^*$ must be zero,
and the bracket must contain zero.

In order to compute the indeterminacy of 
$\langle \theta_4, \eta^2, \theta_4 \rangle$, we must consider
the product of $\theta_4$ with elements of $\pi_{33,18}$.
There are several cases to consider.

First consider $\{\D h_1^2 h_3\}$.  The product
$\theta_4 \{\D h_1^2 h_3\}$ is detected in Adams filtration at least
$10$, but there are no elements in sufficiently high filtration.

Next consider $\nu \theta_4$ detected by $p$.
The product $\theta_4^2$ is zero \cite{Xu16},
so $\nu \theta_4^2$ is also zero.

Finally, consider $\eta \eta_5$ detected by $h_1^2 h_5$.
Table \ref{tab:Toda} shows that $\langle \eta, 2, \theta_4 \rangle$
detects $\eta_5$.
Shuffle to obtain
\[
\eta \eta_5 \theta_4 = \eta \langle \eta, 2, \theta_4 \rangle \theta_4 =
\eta^2 \langle 2, \theta_4, \theta_4 \rangle.
\]
The bracket $\langle 2, \theta_4, \theta_4 \rangle$ is zero
because it is contained in $\pi_{61,32} = 0$.
\end{proof}

\begin{lemma}
\label{lem:eta^2,theta4,eta^2,theta4}
\revdeg{66, 36}
The Toda bracket $\langle \eta^2, \theta_4, \eta^2, \theta_4 \rangle$
is detected by $\D_1 h_3^2$.
\end{lemma}

\begin{proof}
Table \ref{tab:Massey} shows that 
$\D_1 h_3^2$ equals $\langle h_1^2, h_4^2, h_1^2, h_4^2 \rangle$.
Therefore, $\D_1 h_3^2$ detects 
$\langle \eta^2, \theta_4, \eta^2, \theta_4 \rangle$,
if the Toda bracket is well-defined.

In order to show that the Toda bracket is well-defined, we 
need to know that the 
subbrackets $\langle \eta^2, \theta_4, \eta^2 \rangle$
and $\langle \theta_4, \eta^2, \theta_4 \rangle$ contain zero.
These are handled by
Lemmas \ref{lem:eta^2,theta4,eta^2} and \ref{lem:theta4,eta^2,theta4}.
\end{proof}

\begin{lemma}
\label{lem:teta^2kappabar,8,kappabar2}
\revdeg{67, 36}
The Toda bracket 
$\langle \tau \eta^2 \kappabar, 8, \kappabar_2 \rangle$
contains zero,
and its indeterminacy is generated by multiples of
$\tau \eta^2 \kappabar$.
\end{lemma}

\begin{proof}
The bracket
$\langle \tau \eta^2 \kappabar, 8, \kappabar_2 \rangle$
contains
$\tau \eta \kappabar \langle \eta, 2, 4 \kappabar_2 \rangle$.
Lemma \ref{lem:eta,2,4kappabar2} shows that
this expression contains zero.

It remains to show that $\kappabar_2 \cdot \pi_{23, 12}$ equals zero.
There are several cases to consider.

First, the product $\tau \sigma \eta_4 \kappabar_2$ in $\pi_{60,32}$ 
could only be detected by $\tau^4 g^3$ or $\tau^2 d_0^2 l$.
Comparison to $\tmf$ rules out both possibilities.
Therefore,
$\tau \sigma \eta_4 \kappabar_2$ is zero.

Second, the product $\kappabar \kappabar_2$ in $\pi_{64,35}$ must be
detected in filtration at least $9$, since $\tau g g_2$ equals zero,
so it could only be detected by $h_1^2 (\D e_1 + C_0)$.
This implies that $\tau \nu \kappabar \kappabar_2$ is zero.

Third, we must consider the product $\rho_{23} \kappabar_2$.
Table \ref{tab:Toda} shows that
the Toda bracket $\langle \sigma, 16, 2 \rho_{15} \rangle$
detects $\rho_{23}$.
Then
$\rho_{23} \kappabar_2$
is contained in
\[
\langle \sigma, 16, 2 \rho_{15} \rangle \kappabar_2 =
\sigma \langle 16, 2 \rho_{15}, \kappabar_2 \rangle.
\]
The latter bracket is contained in $\pi_{60,32}$.
As above, comparison to $\tmf$ shows that the expression
is zero.
\end{proof}

\begin{lemma}
\label{lem:eta,nu,ttheta4.5kappabar}
\revdeg{70, 37}
The Toda bracket $\langle \eta, \nu, \tau \theta_{4.5} \kappabar \rangle$
is detected by $\tau h_1 D'_3$.
\end{lemma}

\begin{proof}
Table \ref{tab:Adams-d3}
shows that $d_3(\tau D'_3)$
equals
$\tau^2 M h_2 g$.
The Moss Convergence Theorem \ref{thm:Moss} implies that
$\tau h_1 D'_3$ detects the Toda bracket.
\end{proof}

\begin{lemma}
\label{lem:eta,nu,t^2h2C'}
\revdeg{71, 37}
There exists an element $\alpha$ in $\pi_{66,34}$
detected by $\tau^2 h_2 C'$ such that
the Toda bracket $\langle \eta, \nu, \alpha \rangle$
is defined and detected by $\tau h_1 p_1$.
\end{lemma}

\begin{proof}
The differential $d_5(\tau p_1 + h_0^2 h_3 h_6) = \tau^2 h_2^2 C'$
and the Moss Convergence Theorem \ref{thm:Moss} 
establish that the Toda bracket is detected by $\tau h_1 p_1$, provided
that the Toda bracket is well-defined.

Let $\alpha$ be an element of $\pi_{66,34}$ that is detected
by $\tau^2 h_2 C'$.  Then $\nu \alpha$ does not necessarily 
equal zero; it could be detected in higher filtration by
$\tau^2 h_2 B_5 + h_2 D'_2$.
Then we can adjust our choice of $\alpha$ by an element detected
by $\tau^2 B_5 + D'_2$ to ensure that $\nu \alpha$ is zero.
\end{proof}

\begin{lemma}
\label{lem:sigma^2,2,T,tkappabar}
\revdeg{72, 38}
The Toda bracket
$\langle \sigma^2, 2, \{t\}, \tau \kappabar \rangle$
is detected by $h_4 Q_2 + h_3^2 D_2$.
\end{lemma}

\begin{proof}
The subbracket $\langle \sigma^2, 2, \{ t \} \rangle$
contains zero by comparison to $C\tau$, and its indeterminacy
is generated by $\sigma^3 \theta_4 = 4 \sigma \kappabar_2$
detected by $\tau g n$.
The subbracket $\langle 2, \{t\}, \tau \kappabar \rangle$
is strictly zero because it cannot be detected by
$h_0 h_2 h_5 i$ by comparison to $\tmf$.
This shows that the desired four-fold Toda bracket is well-defined.

Consider the relation
\[
\eta \langle \sigma^2, 2, \{t\}, \tau \kappabar \rangle \subseteq
\langle \langle \eta, \sigma^2, 2 \rangle, \{t\}, \tau \kappabar \rangle.
\]
Let $\alpha$ be any element of $\langle \eta, \sigma^2, 2 \rangle$.
Table \ref{tab:Toda} shows that $\alpha$ 
is detected by $h_1 h_4$ and 
equals either $\eta_4$ or $\eta_4 + \eta \rho_{15}$.
By inspection, the indeterminacy of
$\langle \alpha, \{t\}, \tau \kappabar \rangle$
equals $\tau \kappabar \cdot \pi_{53,29}$, which 
is detected in Adams filtration at least 14.  (In fact, 
the indeterminacy is non-zero, since it contains both
$\tau \kappabar \cdot \{M c_0\}$ detected by
$\tau M d_0^2$ and also
$\tau \kappabar \cdot \{\D h_1 d_0^2\}$ detected by
$\tau^2 \D h_1 d_0 e_0^2$.)

Table \ref{tab:Toda} shows that
$\langle \alpha, \{t\}, \tau \kappabar \rangle$ is detected
by $h_1 h_4 Q_2$.
Together with the partial analysis of the indeterminacy in the
previous paragraph, this shows that 
$\langle \alpha, \{t\}, \tau \kappabar \rangle$ does
not contain zero.

Then 
$\eta \langle \sigma^2, 2, \{t\}, \tau \kappabar \rangle$
also does not contain zero, and the only possibility is that
$\langle \sigma^2, 2, \{t\}, \tau \kappabar \rangle$ 
is detected by $h_4 Q_2 + h_3^2 D_2$.
\end{proof}

\begin{lemma}
\label{lem:theta4,theta4,kappa}
\revdeg{75, 40}
The Toda bracket $\langle \theta_4, \theta_4, \kappa \rangle$
equals zero.
\end{lemma}

\begin{proof}
The Massey product $\langle h_4^2, h_4^2, d_0 \rangle$ equals zero,
since
\[
h_1^2 \langle h_4^2, h_4^2, d_0 \rangle =
\langle h_1^2, h_4^2, h_4^2 \rangle d_0 = 0,
\]
while $h_1^2 x_{75,7}$ is not zero.
The Moss Convergence Theorem \ref{thm:Moss} then implies
that $\langle \theta_4, \theta_4, \kappa \rangle$
is detected in Adams filtration at least $8$.

The only element in sufficiently high filtration is $P h_1^4 h_6$.
However,
\[
\eta^2 \langle \theta_4, \theta_4, \kappa \rangle =
\langle \eta^2, \theta_4, \theta_4 \rangle \kappa = 0,
\]
while $h_1^2 \cdot P h_1^4 h_6$ is not zero.
Then 
$\langle \theta_4, \theta_4, \kappa \rangle$ must contain zero because
there are no remaining possibilities.

The indeterminacy can be computed by inspection, using that
$\theta_4 \theta_{4.5}$ is zero by comparison to $C\tau$.
\end{proof}

\begin{lemma}
\label{lem:kappa,2,theta5}
\revdeg{77, 40}
The Toda bracket
$\langle \kappa, 2, \theta_5 \rangle$
is detected by $h_6 d_0$.
\end{lemma}

\begin{proof}
The differential $d_3(h_0 h_4) = h_0 d_0$ implies that
$\langle \kappa, 2, \theta_5 \rangle$
is detected by $h_0 h_4 \cdot h_5^2 = 0$ in filtration 4.
In other words, the Toda bracket is detected in Adams filtration
at least 5.

The element $h_1 h_6 d_0$ detects
$\langle \eta \kappa, 2, \theta_5 \rangle$, 
using the Adams differential $d_2(h_6) = h_0 h_5^2$.
This expression contains
$\eta \langle \kappa, 2, \theta_5 \rangle$, which shows that
$\langle \kappa, 2, \theta_5 \rangle$ is detected 
in filtration at most 5.

The only possibility is that the Toda bracket is detected by $h_6 d_0$. 
\end{proof}

\rev{
\begin{lemma}
\label{lem:tm1,eta,2}
\revdeg{79, 42}
There exists an element $\mu$ in $\pi_{77,41}$ that is detected by
$\tau m_1$ such that $\eta \mu$ is zero and $\mu$ is not divisible by $\tau$.
Moreover, the Toda bracket $\langle \mu, \eta, 2 \rangle$ 
contains zero or is detected by $\tau^2 M e_0^2$, and
its indeterminacy is detected by $h_0 h_2 x_{76,6}$.
\end{lemma}

\begin{proof}
Let $\mu'$ be an element of $\pi_{77,42}$ that is detected by
$m_1$.  Then $\tau \mu'$ is detected by $\tau m_1$, and 
$\eta \mu'$ is detected by $h_1 m_1$.
Table \ref{tab:tau-extn} shows that there is a hidden
$\tau$ extension from $h_1 m_1$ to $M \D h_1^2 h_3$.
Therefore, $\tau \eta \mu'$ is detected by $M \D h_1^2 h_3$.

Now let $\mu''$ be an element of $\pi_{77,41}$ that is detected by
$M \D h_1 h_3$.  Then $\eta \mu''$ is also detected by $M \D h_1^2 h_3$.
This shows that $\eta (\tau \mu' + \mu'')$ is zero because there 
are no possible detecting elements in higher filtration.

Choose $\mu$ to be $\tau \mu' + \mu''$.
Note that $\mu''$ is not divisible by $\tau$ because
inclusion of the bottom cell of $C\tau$ takes $M \D h_1 h_3$ to
a non-zero element.  Therefore, $\mu$ is also not divisible
by $\tau$.

Now that $\mu$ is defined, it remains to study the Toda bracket.
We begin with an analysis of its indeterminacy, which is generated
by $\tau \eta^2 \cdot \mu$ and the multiples of $2$ in 
$\pi_{79,42}$.  The first expression is zero by the construction
of $\mu$. 
Let $\alpha$ be an element of $\pi_{79,42}$ that is detected by
$h_2 x_{76,6}$, so $2 \alpha$ is detected by $h_0 h_2 x_{76,6}$.
Tables \ref{tab:2-extn} and \ref{tab:2-extn-possible}
show that there are no hidden $2$ extensions in the $79$-stem with
weight $42$.  Therefore, the indeterminacy is generated by $2 \alpha$.



Inclusion of the bottom cell of $C\tau$ takes the bracket to
$\langle M \D h_1 h_3, h_1, h_0 \rangle$.
Machine-generated data \cite{Wang19} shows that this bracket 
equals $\{0, h_0 h_2 x_{76,6} \}$ 
in $C\tau$.  

Let $\beta$ be any element of $\langle \mu, \eta, 2 \rangle$.
It is possible that $\beta$ maps to $h_0 h_2 x_{76,6}$ under
inclusion of the bottom cell of $C\tau$.  In that case,
$\beta + 2 \alpha$ also belongs to $\langle \mu, \eta, 2 \rangle$
and must map to zero under inclusion of the bottom cell of $C\tau$.

In either case, the original Toda bracket
contains an element 
that maps to zero under inclusion of the bottom cell of $C\tau$,
and that element is therefore divisible by $\tau$.
By inspection, the only possible detecting elements are
$\tau^2 M e_0^2$ and $\tau^3 \D h_1 e_0^2 g$.
The latter option is ruled out by comparison to $\mmf$.
\end{proof}
}

\begin{lemma}
\label{lem:2,eta,th1^2x76,6}
\revdeg{80, 42}
The Toda bracket
$\langle 2, \eta, \tau \eta \{h_1 x_{76,6} \} \rangle$ is detected by
$\tau h_1 x_1$.
\end{lemma}

\begin{proof}
Let $\alpha$ be an element of $\pi_{77,41}$ that is detected by
$h_1 x_{76,6}$.
First we must show that the Toda bracket is well-defined.

Note that $2 \alpha$ is zero because there are no 
$2$ extensions in $\pi_{77,41}$ in sufficiently high
Adams filtration.
Now consider the shuffle
\[
\tau \eta^2 \alpha = \langle 2, \eta, 2 \rangle \alpha =
2 \langle \eta, 2, \alpha \rangle.
\]
Table \ref{tab:Toda} shows that $\langle \eta, 2, \alpha \rangle$
is detected by $h_0 h_2 x_{76,6}$, but this element does
not support a hidden $2$ extension.  This shows that
$\tau \eta^2 \alpha$ is zero and that the Toda bracket
is well-defined.

Finally, use the Adams differential $d_4(h_0 e_2) = \tau h_1^3 x_{76,6}$
and the relation $h_0 \cdot h_0 e_2 = \tau h_1 x_1$ to compute
the Toda bracket.
\end{proof}

\rev{
\begin{lemma}
\label{lem:2,eta,h2x76,6}
\revdeg{81, 43}
There exists an element $\alpha$ in $\pi_{79,42}$ that is detected by
$h_2 x_{76,6}$ such that
$\eta \alpha$ is zero.  Moreover, the Toda bracket
$\langle 2, \eta, \alpha \rangle$ is zero, with no indeterminacy.
\end{lemma}

\begin{proof}
There is no hidden $\eta$ extension on $h_2 x_{76,6}$ 
because the possible targets $\tau^3 d_0 e_0^2 l$ and
$\tau^5 g^4$ are ruled out by comparison to $\mmf$.
Therefore, $\alpha$ exists.

The Massey product $\langle h_0, h_1, h_2 x_{76,6} \rangle$
has no indeterminacy by inspection.  Consequently,
\[
\langle h_0, h_1, h_2 x_{76,6} \rangle =
\langle h_0, h_1, h_2 \rangle x_{76,6} = 0.
\]
The Moss Convergence Theorem \ref{thm:Moss} implies that
the Toda bracket $\langle 2, \eta, \alpha \rangle$ is detected
in filtration $9$ or greater.
The possible detecting elements are $P h_1^2 h_6 c_0$ and
$\D^2 h_1 d_1$.  In either of these cases, the Toda bracket
would be detected by inclusion of the bottom cell of $C\tau$,
but the corresponding bracket is zero in $C\tau$.

The indeterminacy is generated by $\tau \eta^2 \alpha$
and the multiples of $2$ in $\pi_{81,43}$.  The first expression
is zero by the choice of $\alpha$.  Tables \ref{tab:2-extn}
and \ref{tab:2-extn-possible} show that there are no
multiples of $2$ in $\pi_{81,43}$.
\end{proof}
}


\begin{lemma}
\label{lem:2,sigma^2,th2^2C'}
\revdeg{84, 45}
The Toda bracket $\langle 2, \sigma^2, \{\tau h_2^2 C'\} \rangle$
equals zero, with no indeterminacy.
\end{lemma}

\begin{proof}
Let $\alpha$ be an element of $\pi_{66,35}$ that is detected by
$\tau h_2 C'$, so $\nu \alpha$ is the unique element
that is detected by $\tau h_2^2 C'$.
We consider the Toda bracket
$\langle 2, \sigma^2, \nu \alpha \rangle$.
By inspection, the indeterminacy is zero, so
the bracket equals
$\langle 2, \sigma^2, \nu \rangle \alpha$, which equals
$\langle \alpha, 2, \sigma^2 \rangle \nu$.

Apply the Moss Convergence Theorem \ref{thm:Moss}
with the Adams $d_2$ differential to see
that the Toda bracket $\langle \alpha, 2, \sigma^2 \rangle$
is detected by $0$ in Adams filtration $9$, but it could be detected
by a non-zero element in higher filtration.  However, this shows
that $\langle \alpha, 2, \sigma^2 \rangle \nu$ is zero by inspection.
\end{proof}

\begin{lemma}
\label{lem:2,sigma^2,h3(De1+C0)}
\revdeg{84, 45}
The Toda bracket $\langle 2, \sigma^2, \{h_3(\D e_1 + C_0)\} \rangle$
equals zero, with no indeterminacy.
\end{lemma}

\begin{proof}
Let $\beta$ be an element of $\pi_{62,33}$ that is detected by
$\D e_1 + C_0$, so $\sigma \beta$ is the unique element that is
detected by $h_3(\D e_1 + C_0)$.
We consider the Toda bracket
$\langle 2, \sigma^2, \sigma \beta \rangle$.
By inspection, the indeterminacy is zero, so
the bracket equals
$\langle 2, \sigma^2, \beta \rangle \sigma$.

Apply the Moss Convergence Theorem \ref{thm:Moss} 
with the Adams $d_2$ differential to see that the Toda bracket
$\langle 2, \sigma^2, \beta \rangle$ is detected by $0$ in 
Adams filtration $9$, but it could be detected by a non-zero
element in higher filtration. 
Then the only possible non-zero value for
$\langle 2, \sigma^2, \beta \rangle \sigma$ is
$\{M \D h_1 h_3\} \sigma$.
Table \ref{tab:misc-extn} shows that $M \D h_1 h_3$ detects
$\{\D h_1 h_3\} \theta_{4.5}$, so
$\sigma \{M \D h_1 h_3\}$ equals $\sigma \{\D h_1 h_3\} \theta_{4.5}$,
which equals zero.
\end{proof}

\begin{lemma}
\label{lem:tetakappabar^2,2,4kappabar2}
\revdeg{86, 46}
The Toda bracket $\langle \tau \eta \kappabar^2, 2, 4 \kappabar_2 \rangle$
is detected by $M \D h_0^2 e_0$.
\end{lemma}

\begin{proof}
Table \ref{tab:Massey} shows that the Massey product
$\langle \D h_0^2 e_0, h_0^2, h_0 g_2 \rangle$ equals the element
$M \D h_0^2 e_0$.
Now apply the Moss Convergence Theorem \ref{thm:Moss}, using that
Table \ref{tab:eta-extn} shows that
$\D h_0^2 e_0$ detects $\tau \eta \kappabar^2$.
\end{proof}

\begin{lemma}
\label{lem:th0Q3,nu4,eta}
\revdeg{87, 46}
There exists an element $\alpha$ in $\pi_{67,36}$ that is detected
by $h_0 Q_3 + h_0 n_1$ such that 
$h_1^2 c_3$ detects the Toda bracket
$\langle \tau \alpha, \nu_4, \eta \rangle$.
\end{lemma}

\begin{proof}
A consequence of the proof of Lemma \ref{lem:d4-h1c3} is that
there exists $\alpha$ in $\pi_{67,36}$ that is detected by
$h_0 Q_3 + h_0 n_1$ such that
the product $\tau \nu_4 \alpha$ is zero.
Therefore,
$h_1^2 c_3$ detects the Toda bracket
$\langle \tau \alpha, \nu_4, \eta \rangle$ because
of the Adams differential $d_4(h_1 c_3) = \tau h_0 h_2 h_4 Q_3$.
\end{proof}
